\subsection{Docker Compose}
Primeramente deberemos crear un archivo \texttt{docker-compose.yml} con el siguiente contenido:
\begin{verbatim}
services:
  samba-ad-dc:
    image: vanleuken/samba4-ad-dc
    container_name: samba-ad-dc
    hostname: ad
    environment:
      - DOMAIN=example.local
      - REALM=EXAMPLE.LOCAL
      - ADMIN_PASSWORD=Passw0rd!
    ports:
      - "53:53/udp"
      - "53:53/tcp"
      - "88:88"
      - "135:135"
      - "137-138:137-138/udp"
      - "139:139"
      - "389:389"
      - "445:445"
      - "464:464"
      - "636:636"
      - "3268-3269:3268-3269"
      - "49152-49155:49152-49155"
    volumes:
      - samba_data:/var/lib/samba
      - samba_etc:/etc/samba

  lam:
    image: ghcr.io/ldapaccountmanager/lam:latest
    container_name: ldap-account-manager
    ports:
      - "8080:80"
    environment:
      - LDAP_SERVER=samba-ad-dc
      - LDAP_BASE_DN=DC=example,DC=local
    depends_on:
      - samba-ad-dc

volumes:
  samba_data:
  samba_etc:
\end{verbatim}
Este archivo define dos servicios: \texttt{samba-ad-dc} para el servidor Samba
Active Directory y \texttt{lam} para la interfaz web de administración LDAP Account Manager.
\subsection{Iniciar los contenedores}
Para iniciar los contenedores, ejecutamos el siguiente comando en el directorio donde se encuentra
el archivo \texttt{docker-compose.yml}:
\begin{verbatim}
docker-compose up -d
\end{verbatim}
Esto descargará las imágenes necesarias y levantará los contenedores en segundo plano.
\subsection{Acceder a LDAP Account Manager}
Una vez que los contenedores estén en funcionamiento, podemos acceder a la interfaz web de LDAP
Account Manager abriendo un navegador web y navegando a \url{http://localhost:808
0}. El nombre de usuario predeterminado es \texttt{lam} y la
contraseña es \texttt{lam}.
\subsection{Configuración inicial}
Al iniciar sesión por primera vez, se nos pedirá configurar la conexión LDAP. Utilizamos los
siguientes parámetros:
\begin{itemize}
    \item \textbf{Servidor LDAP:} \texttt{samba-ad-dc}
    \item \textbf{Base DN:} \texttt{DC=example,DC=local}
    \item \textbf{Usuario:} \texttt{Administrator@example.local}
    \item \textbf{Contraseña:} \texttt{Passw0rd!}
\end{itemize}
Después de configurar la conexión, podremos gestionar usuarios, grupos y otros objetos en el
servidor Samba Active Directory a través de la interfaz web de LDAP Account Manager.

